% ! TEX root = ../m3-19-20.tex

\chapter{Transformata Laplace și transformata Z}

\section{Transformata Laplace}

\subsection{Definiții și proprietăți}

Transformata Laplace este o transformare integrală, care se poate aplica unor
funcții speciale, numite \emph{funcții original}.

\begin{definition}\label{def:orig}\index{funcție!original}
  O funcție $ f : \RR \to \CC $ se numește \emph{original} dacă:
  \begin{enumerate}[(a)]
  \item $ f(t) = 0 $ pentru orice $ t < 0 $;
  \item $ f $ este continuă (eventual pe porțiuni) pe intervalul $ [0, \infty) $;
  \item Este mărginită de o exponențială, adică există $ M > 0 $ și $ s_0 \geq 0 $
    astfel încît:
    \[
      |f(t)| \leq Me^{s_0 t}, \quad \forall t \geq 0.
    \]

    Numărul $ s_0 $ se mai numește \emph{indicele de creștere} al
    funcției $ f $.

    Vom nota cu $ \kal{O} $ mulțimea funcțiilor original.
  \end{enumerate}
\end{definition}

Pornind cu o funcție original, definiția transformatei Laplace este:
\begin{definition}\label{def:lapl}\index{transformata!Laplace}
  Păstrînd contextul și notațiile de mai sus, fie $ f \in \kal{O} $ și
  mulțimea:
  \[
    S(s_0) = \{ s \in \CC \mid \dr{Re}(s) > s_0 \}.
  \]

  Funcția:
  \[
    F : S(s_0) \to \CC, \quad F(s) = \int_0^\infty f(t) e^{-st} dt
  \]
  se numește \emph{transformata Laplace} a lui $ f $ sau \emph{imaginea Laplace}
  a originalului $ f $.

  Vom mai folosi notația $ F = \kal{L} f $ sau, explicit, $ \kal{L} f(t) = F(s). $
\end{definition}

Proprietățile esențiale ale transformatei Laplace sînt date mai jos. Fiecare dintre ele
va fi folosită pentru a calcula o transformată Laplace pentru o funcție care nu se
regăsește direct într-un tabel de valori.

\begin{itemize}
\item \textbf{Liniaritate:} $ \kal{L}(\alpha f + \beta g) = \alpha \kal{L} f + \beta \kal{L} g $,
  pentru orice $ \alpha, \beta \in \CC $, iar $ f, g $ funcții original;
\item \textbf{Teorema asemănării:}
  \[
    \kal{L} f(\alpha t) = \dfrac{1}{\alpha} F \Big( \dfrac{s}{\alpha} \Big);
  \]
\item \textbf{Teorema deplasării:}
  \[
    \kal{L} \big( f(t) e^{s_0t} \big) = F(s - s_0);
  \]
\item \textbf{Teorema întîrzierii:} Definim \emph{întîrziata cu $ \tau $ } a funcției
  $ f \in \kal{O} $ prin:
  \[
    f_\tau(t) =
    \begin{cases}
      0, & t < \tau \\
      f(t - \tau), & t \geq \tau
    \end{cases}.
  \]
  Atunci, dacă $ \Lap f(t) = F(s) $, $ \Lap f_\tau (t) = e^{-st} F(s) $;
\item \textbf{Teorema derivării imaginii:}
  \[
    \Lap \big( t^n f(t) \big) = (-1)^n F^{(n)} (s);
  \]
\item \textbf{Teorema integrării originalului:} Fie $ f \in \kal{O}, \Lap f(t) = F(s) $ și
  $ g(t) = \ds\int_0^t f(\tau) d\tau $. Atunci:
  \[
    \Lap g(t) = \dfrac{1}{s} F(s);
  \]
\item \textbf{Teorema integrării imaginii:} Fie $ \Lap f(t) = F(s) $ și $ G $ o primitivă a
  lui $ F $ în $ S(s_0) $, cu $ G(\infty) = 0 $. Atunci:
  \[
    \Lap \dfrac{f(t)}{t} = -G(s).
  \]
\end{itemize}

%%%%%%%%%%%%%%%%%%%%%%%%%%%%%%%%%%%%%%%%%%%%%%%%%%%%%%%%%%%%%%%%%%%%%% 
\newpage
\subsection{Tabel de transformate Laplace}

\indent\indent În tabelul de mai jos, vom considera funcțiile $ f(t) $ ca fiind funcții original,
adică nule pentru argument negativ. Echivalent, putem gîndi $ f(t) $ ca fiind, de fapt,
înmulțite cu \emph{funcția lui Heaviside}:
\[
  u(t) =
  \begin{cases}
    1, & t \geq 0 \\
    0, & t < 0 
  \end{cases}.
\]

\begin{center}
\begin{tabular}{c | c}
  $ f(t) = \Lap^{-1} (F(s)) $ & $ F(s) = \Lap(f(t)) $ \\
  \hline\hline
  $ u(t) $ & $ \dfrac{1}{s} $ \\
  $ u(t - \tau) $ & $ \dfrac{1}{s} e^{-\tau s} $ \\
  $ t $ & $ \dfrac{1}{s^2} $ \\
  $ t^n $ & $ \dfrac{n!}{s^{n+1}} $ \\
  $ t^n e^{-\alpha t} $ & $ \dfrac{n!}{(s + \alpha)^{n+1}} $ \\
  $ e^{-\alpha t} $ & $ \dfrac{1}{s + \alpha} $ \\
  $ \sin (\omega t) $ & $ \dfrac{\omega}{s^2 + \omega^2} $ \\
  $ \cos (\omega t) $ & $ \dfrac{s}{s^2 + \omega^2} $ \\
  $ \sinh(\alpha t) $ & $ \dfrac{\alpha}{s^2 - \alpha^2} $ \\
  $ \cosh(\alpha t) $ & $ \dfrac{s}{s^2 - \alpha^2} $ \\
  $ e^{-\alpha t}\sin(\omega t) $ & $ \dfrac{\omega}{(s + \alpha)^2 + \omega^2} $ \\
  $ e^{-\alpha t}\cos(\omega t) $ & $ \dfrac{s + \alpha}{(s + \alpha^2) + \omega^2} $ \\
  $ \ln t $ & $ -\dfrac{1}{s} \big( \ln s + \gamma\footnotemark \big) $
\end{tabular}
\end{center}

\footnotetext{Constanta Euler-Mascheroni, $ \gamma \simeq 0,577\dots \in \RR - \QQ $}


%%%%%%%%%%%%%%%%%%%%%%%%%%%%%%%%%%%%%%%%%%%%%%%%%%%%%%%%%%%%%%%%%%%%%% 
\newpage
\subsection{Exerciții}

1. Calculați transformatele Laplace pentru funcțiile (presupuse original):
\begin{enumerate}[(a)]
\item $ f(t) = 1, t \geq 0 $;
\item $ f(t) = t, t \geq 0 $;
\item $ f(t) = t^n, n \in \NN $;
\item $ f(t) = e^{at}, t \geq 0, a \in \RR $;
\item $ f(t) = \sin(at), t \geq 0, a \in \RR $.
\end{enumerate}

\emph{Soluție:} (a) Avem direct din definiție:
\[
  F(s) = \int_0^\infty e^{-st} dt = \lim_{n \to \infty} \int_0^n e^{-st} dt = \frac{1}{s}, s > 0.
\]

(b) Integrăm prin părți și obținem $ F(s) = \dfrac{1}{s^2} $.

(c) Facem substituția $ st = \tau $ și găsim:
\[
  F(s) = \int_0^\infty e^{-st} t^n dt = \frac{1}{s^{n+1}} \int_0^\infty e^{-\tau} \tau^n d\tau = %
  \frac{n!}{s^{n+1}},
\]
pentru $ s > 0 $, folosind funcția Gamma a lui Euler:
\[
  \Gamma(a) = \int_0^\infty x^{a-1}e^{-x} dx, a > 0, \quad \Gamma(n + 1) = n!, n \in \NN.
\]

(d) $ F(s) = \ds\int_0^\infty e^{at} e^{-st} dt = \ds\int_0^\infty e^{-(s - a)t} dt = %
\dfrac{1}{s-a}, s > a $.

(e) Integrăm prin părți și ajungem la:
\[
  F(s) = \frac{1}{a} - \frac{s^2}{a^2} F(s) \Rightarrow F(s) = \frac{a}{s^2 + a^2}, s > 0.
\]

\vspace{1cm}

2. Folosind tabelul de valori și proprietățile, să se determine transformatele Laplace
pentru funcțiile (presupuse original):
\begin{enumerate}[(a)]
\item $ f(t) = 5 $;
\item $ f(t) = 3t + 6t^2 $;
\item $ f(t) = e^{-3t} $;
\item $ f(t) = 5e^{-3t} $;
\item $ f(t) = \cos (5t) $;
\item $ f(t) = \sin (3t) $;
\item $ f(t) = 3(t-1) + e^{-t-1} $;
\item $ f(t) = 3t^3(t - 1) + e^{-5t} $;
\item $ f(t) = 5e^{-3t} \cos (5t) $;
\item $ f(t) = e^{2t} \sin (3t) $;
\item $ f(t) = te^{-t} \cos (4t) $;
\item $ f(t) = t^2 \sin (3t) $;
\item $ f(t) = t^3 \cos t $.
\end{enumerate}

\emph{Indicații:} În majoritatea cazurilor, se folosește tabelul și proprietatea de liniaritate.
În plus:

(i, j) Folosim $ \Lap (e^{at} f(t)) = F(s - a) $;

(k) Folosim $ \Lap (t f(t)) = -\dfrac{d}{ds} \Lap(f(t)) $;

(l) Folosim $ \Lap(t^n f(t)) = (-1)^n \dfrac{d^n F(s)}{ds^n} $;

\vspace{1cm}

3. Folosind teorema derivării imaginii, să se determine transformatele Laplace pentru funcțiile
(presupuse original):
\begin{enumerate}[(a)]
\item $ f(t) = t $;
\item $ f(t) = t^2 $;
\item $ f(t) = t\sin t $;
\item $ f(t) = te^t $.
\end{enumerate}

\emph{Indicație:} Conform proprietății de derivare a imaginii, avem:
\[
  \Lap (t f(t)) = -\frac{d}{ds} \Lap (f(t)).
\]

\vspace{1cm}

4. Folosind teorema integrării originalului, să se determine transformatele Laplace pentru
funcțiile (presupuse original):
\begin{enumerate}[(a)]
\item $ f(t) = \ds\int_0^t \cos (2\tau) d\tau $;
\item $ f(t) = \ds\int_0^t e^{3 \tau} \cos(2 \tau) d\tau $;
\item $ f(t) = \ds\int_0^t \tau e^{-3\tau} d\tau $.
\end{enumerate}

\emph{Indicație:} Conform proprietății de integrare a originalului,
avem:
\[
  \Lap \int_0^t f(\tau) d\tau = \frac{F(s)}{s}.
\]


%%%%%%%%%%%%%%%%%%%%%%%%%%%%%%%%%%%%%%%%%%%%%%%%%%%%%%%%%%%%%%%%%%%%%% 

\subsection{Aplicații ale transformatei Laplace}

Principala aplicație a transformatei Laplace este pentru rezolvarea ecuațiilor și
sistemelor diferențiale de ordinul întîi sau superior.

Aceste aplicații se bazează pe calculele care se pot obține imediat din definiția
transformatei Laplace și a proprietăților sale:
\begin{align*}
  \Lap f' &= s \Lap f - f(0) \\
  \Lap f^{''} &= s^2 \Lap f - sf(0) - f'(0).
\end{align*}

De fapt, în general, avem:
\[
  \Lap f^{(n)} = s^n \Lap f - s^{n-1} f(0) - s^{n-2} f'(0) - \dots - f^{(n-1)}(0).
\]

De asemenea, pentru integrale, știm deja teorema integrării originalului:
\[
  \Lap \int_0^t f(\tau) d\tau = \dfrac{1}{s} F(s), \quad F(s) = \Lap f(t).
\]

Rezultă, folosind transformarea inversă:
\[
  \int_0^t f(\tau) d\tau = \Lap^{-1} \Big( \dfrac{1}{s} F(s) \Big).
\]

De exemplu, pentru a rezolva ecuația diferențială:
\[
  y^{\prime \prime} + ay^\prime + by = r(t), \quad y(0) = K_0, y'(0) = K_1,
\]
aplicăm transformata Laplace și folosim proprietățile de mai sus. Fie $ Y = \Lap y(t) $

Se obține ecuația algebrică:
\[
  (s^2 Y - sy(0) - y'(0)) + a(sY - y(0)) + bY = R(s),
\]
unde $ R(s) = \Lap r $. Forma echivalentă este:
\[
  (s^2 + as + b) Y = (s + a)y(0) + y'(0) + R(s).
\]

Împărțim prin $ s^2 + as + b $ și folosim formula:
\[
  Q(s) = \frac{1}{s^2 + as + b} =
  \frac{1}{(s + \frac{1}{2}a)^2 + b - \frac{1}{4}a^2},
\]
de unde rezultă:
\[
  Y(s) = \big( (s + a) y(0) + y'(0)\big) Q(s) + R(s) Q(s).
\]

În forma aceasta, descompunem $ Y(s) $ în fracții simple, dacă este nevoie și
folosim tabelul de transformate Laplace, pentru a afla $ y = \Lap^{-1}(Y) $.

\vspace{1cm}

De exemplu:
\[
  y^{\prime\prime} - y = t, \quad y(0) = 1, y'(0) = 1.
\]

\emph{Soluție:} Aplicăm transformata Laplace și ajungem la ecuația:
\begin{align*}
  s^2 Y - sy(0) - y'(0) - Y &= \frac{1}{s^2} \\
  (s^2 - 1)Y &= s + 1 + \frac{1}{s^2}.
\end{align*}

Rezultă $ Q = \dfrac{1}{s^2 - 1} $ și ecuația devine:
\begin{align*}
  Y &= (s + 1)Q + \frac{1}{s^2} Q \\
    &= \frac{s + 1}{s^2 - 1} + \frac{1}{s^2(s^2 - 1)} \\
    &= \frac{1}{s - 1} + \Big( \frac{1}{s^2 - 1} - \frac{1}{s^2} \Big)
\end{align*}

Folosind tabelul și proprietățile transformatei Laplace, obținem soluția:
\begin{align*}
  y(t) &= \Lap^{-1} Y \\
       &= \Lap^{-1} \Big( \frac{1}{s^{-1}} \Big) + \Lap^{-1} \Big( \frac{1}{s^2 - 1} \Big) - %
         \Lap^{-1} \Big( \frac{1}{s^2} \Big) \\
       &= e^t + \sinh t - t.
\end{align*}

\vspace{1cm}

%%%%%%%%%%%%%%%%%%%%%%%%%%%%%%%%%%%%%%%%%%%%%%%%%%%%%%%%%%%%%%%%%%%%%% 

\subsection{Exerciții}

\indent\indent 1. Să se rezolve următoarele probleme Cauchy, folosind transformata
Laplace:
\begin{enumerate}[(a)]
\item $ y'(t) + 2y(t) = 4t, y(0) = 1 $ ;
\item $ y'(t) + y(t) = \sin 4t, y(0) = 0 $;
\item $ y^{\prime\prime}(t) + 2y'(t) + 5y(t) = 0, y(0) = 0, y'(0) = 0 $;
\item $ 2y^{\prime\prime}(t) - 6y'(t) + 4y(t) = 3e^{3t}, y(0) = 1, y'(0) = -1 $;
\item $ y^{\prime\prime}(t) - 2y(t) + y(t) = e^t, y(0) = -2, y'(0) = -3 $;
\item $ y^{\prime\prime}(t) + 4y(t) = 3\cos^2(t), y(0) = 1, y'(0) = 2 $.  
\end{enumerate}


\vspace{1cm}

2. Să se rezolve următoarele sisteme diferențiale:
\begin{enumerate}[(a)]
\item $
  \begin{cases}
    x' + x + 4y &= 10 \\
    x - y' - y &= 0 \\
    x(0) = 4, & y(0) = 3
  \end{cases}
  $
\item $
  \begin{cases}
    x' + x - y &= e^t \\
    y' + y - x &= e^t \\
    x(0) = 1, & y(0) = 1
  \end{cases}
  $
\item $
  \begin{cases}
    x' + 2y' + x - y &= 5 \sin t \\
    2x' + 3y' + x - y &= e^t \\
    x(0) = 2, & y(0) = 1
  \end{cases}
  $
\end{enumerate}

\emph{Indicație:} Aplicăm transformata Laplace fiecărei ecuații
și notăm $ \Lap x(t) = X(s) $ și $ \Lap y(t) = Y(s) $. Apoi rezolvăm sistemul
\emph{algebric} obținut cu necunoscutele $ X $ și $ Y $, cărora la final le
aplicăm transformata Laplace inversă.

%%%%%%%%%%%%%%%%%%%%%%%%%%%%%%%%%%%%%%%%%%%%%%%%%%%%%%%%%%%%%%%%%%%%%%
\subsection{Formula de inversare Mellin-Fourier}

În cazul simplu, dacă se dă o imagine Laplace și vrem să recuperăm funcția
original, putem pur și simplu să citim tabelul de transformate de la stînga
la dreapta, eventual după prelucrarea imaginii Laplace, precum am văzut în
exemplele de mai sus. Dar avem nevoie și de o metodă care să funcționeze 
pe cazul general. Mai precis, avem nevoie de o metodă de a \emph{inversa}
transformatele Laplace. Această metodă este dată de formula de inversare
Mellin-Fourier, pe care o prezentăm mai jos într-o variantă simplificată,
mai potrivită pentru aplicații.

\begin{theorem}[Mellin-Fourier]\label{thm:mellin-fourier}
  \index{teorema!Mellin-Fourier}
  Fie $ f $ o funcție original, $ F = \Lap f $.

  Atunci are loc egalitatea:
  \[
    f(t) = \sum_{k=1}^n \Rez\left(F(p)e^{pt}, p_k\right).
  \]
\end{theorem}

Deoarece avem mai multe moduri particulare de a calcula reziduurile, distingem
și aici două calcule speciale:
\begin{itemize}
\item Dacă $ p_1, p_2, \dots, p_r $ sînt poli de ordin $ n_1, n_2, \dots, n_r $, atunci
  formula de inversare se mai scrie:
  \[
    f(t) = \sum_{k=1}^r \dfrac{1}{(n_k - 1)!} \cdot %
    \lim_{p \to p_k} \dfrac{d^{n_k - 1}}{dp^{n_k - 1}} \left[ %
      F(p)(p - p_k)^{n_k} e^{pt} \right];
  \]
\item Dacă imaginea Laplace se poate scrie $ F(p) = \dfrac{A(p)}{B(p)} $,
  unde $ A $ și $ B $ sînt polinoame cu coeficienți reali, astfel încît
  $ \dr{grad}(A) < \dr{grad}B $, atunci formula de inversare devine:
  \[
    f(t) = \sum_k \Rez\left(\dfrac{A(p)}{B(p)}e^{pt}, p_k\right) + %
    \sum_k 2 \dr{Re} \left( \Rez \left( \dfrac{A(p)}{B(p)}e^{pt}, p_k \right) \right),
  \]
  unde prima sumă se calculează pentru polii reali, iar cea de-a doua, pentru
  polii complecși cu partea imaginară pozitivă.
\end{itemize}

\textbf{Exemplu:} Determinați funcția original corespunzătoare imaginii
Laplace:
\[
  F(p) = \dfrac{3p - 4}{p^2 - p - 6}.
\]

\emph{Soluție:} Observăm că avem $ p_1 = 3, p_2 = -2 $ poli simpli.
Atunci rezultă:
\[
  f(t) = \Rez\left( \dfrac{3p - 4}{p^2 - p - 6} e^{pt}, 3\right) + %
  \Rez\left( \dfrac{3p - 4}{p^2 - p - 6} e^{pt}, -2 \right).
\]

Cum ambii sînt poli simpli, reziduurile se calculează folosind
formula din teorma \ref{thm:calc-rez}(3) și se obține în final
$ f(t) = e^{3t} + 2e^{-2t} $.

\begin{remark}\label{rk:no-mellin}
  Această problemă, ca și altele, pot fi rezolvate și fără formula de inversare.
  Trebuie doar atenție și suficient exercițiu astfel încît, în afară de citirea
  \qq{invers} a tabelului de transformate Laplace, să putem folosi și teoreme
  de tip întîrziere, derivarea originalului etc. De exemplu, exercițiul de mai
  sus putea începe cu descompunerea în fracții simple:
  \[
    \dfrac{3p - 4}{p^2 - p - 6} = \dfrac{A}{p - 3} + \dfrac{B}{p + 2}.
  \]
  Se determină $ A $ și $ B $, apoi se folosește tabelul de transformate.
\end{remark}




%%%%%%%%%%%%%%%%%%%%%%%%%%%%%%%%%%%%%%%%%%%%%%%%%%%%%%%%%%%%%%%%%%%%%%

\section{Transformata Z}

Această transformată se definește pe un caz discret, pornind de la:
\begin{definition}\label{def:semnal}\index{semnal!discret}
  Se numește \emph{semnal discret} o funcție $ x : \ZZ \to \CC $ dată
  de $ n \mapsto x_n $ sau, echivalent, $ x(n) $ ori $ x[n] $.
\end{definition}

Mulțimea semnalelor discrete se va nota cu $ S_d $, iar cele cu suport pozitiv
(nule pentru $ n < 0 $ se va nota $ S_d^+ $.

Un semnal particular este \emph{impulsul unitar discret} la momentul $ k $,
definit prin:\index{semnal!unitar}
\[
  \delta_k(n) = \begin{cases} 1, & n = k \\ 0 & n \neq k \end{cases},
\]
definit pentru $ k \in \ZZ $ fixat.

Pentru $ k = 0 $, vom nota $ \delta_0 = \delta $.

\begin{definition}\label{def:intirziat} \index{semnal!intîrziat}
  Fie $ x \in S_d $ și $ k \in \ZZ $ fixat. Semnalul $ y = (x_{n-k}) $ se numește
  \emph{întîrziatul lui $ x $ cu $ k $ momente}.
\end{definition}

O operație foarte importantă, pe care se vor baza unele proprietăți esențiale ale
transformatei $ Z $ este \emph{convoluția:} \index{convoluție}
\begin{definition}\label{def:convolutie}
  Fie $ x, y \in S_d $. Dacă seria $ \ds\sum_{k \in \ZZ} x_{n-k} y_k $
  este convergentă pentru orice $ n \in \ZZ $ și are suma $ z_n $, atunci semnalul
  $ z = (z_n) $ se numește \emph{convoluția} semnalelor $ x $ și $ y $ și
  se notează $ z = x \ast y $.
\end{definition}

Trei proprietăți imediate sînt:
\begin{itemize}
\item $ x \ast y = y \ast x $;
\item $ x \ast \delta = x $;
\item $ (x \ast \delta_k)(n) = x_{n - k} $.
\end{itemize}

Ajungem acum la definiția principală:
\begin{definition}\label{def:z}\index{transformata!Z}
  Fie $ s \in S_d $, cu $ s = (a_n)_n $. Se numește \emph{transformata Z} sau
  \emph{transformata Laplace discretă} a acestui semnal funcția definită prin:
  \[
    L_s(z) = Z_s(z) = \sum_{n \in \ZZ} a_n z^{-n},
  \]
  care se definește în domeniul de convergență al seriei Laurent din definiție.
\end{definition}

Principalele proprietăți pe care le vom folosi în calcule sînt:
\begin{enumerate}[(1)]
\item \textbf{Inversarea transformării $ Z $:} Fie $ s \in S_d^+ $, cu $ s = (a_n) $.
  Presupunem că $ Z_s(z) $ este olomorfă în domeniul $ |z| \in (r, R) $. Atunci
  putem recupera semnalul $ a_n $ prin formula:
  \[
    a_n = \dfrac{1}{2 \pi i} \int_{\gamma} z^{n-1} Z_s(z) dz, \quad n \in \ZZ,
  \]
  unde $ \gamma $ este discul de rază $ \rho \in (r, R) $.
\item \textbf{Teorema de convoluție:} Fie $ s, t \in S_d^+ $. Atunci $ s \ast t = S_d^+ $
  și are loc $ L_{s \ast t} = L_s \cdot L_t $. În particular:
  \[
    L_{s \ast \delta_k}(z) = z^{-k} L_s(z), \quad k \in \ZZ;
  \]
\item \textbf{Prima teoremă de întîrziere:} Pentru $ n\in \NN^\ast $:
  \[
    L_s( z_{t - n}) = z^{-n} L_s(f);
  \]
\item \textbf{A doua teoremă de întîrziere (teorema de deplasare):}
  \[
    L_s(z_{t + n}) = z^n \cdot \Big( L_s(z) - \sum_{t = 0}^{n-1} z_t z^{-t} \Big), \quad n \in \NN^\ast.
  \]
\end{enumerate}

Cîteva transformate uzuale sînt:
\begin{tabular}{c | c}
  $ s $ & $ L_s $ \\
  \hline\hline
  $ \begin{cases}
    h_n = 0, & n < 0 \\
    h_n = 1, & n \geq 0
  \end{cases} $ & $ \dfrac{z}{z - 1} $ \\
        & \\
  $ \delta_k, k \in \ZZ $ & $ \dfrac{1}{z^k} $ \\
        &\\
  $ s = (n)_{n \in \NN} $ & $ \dfrac{z}{(z - 1)^2} $ \\
        &\\
  $ s = (n^2)_{n \in \NN} $ & $ \dfrac{z(z+1)}{(z - 1)^3} $ \\
        &\\
  $ s = (a^n)_{n \in \NN}, a \in \CC $ & $ \dfrac{z}{z - a} $ \\
        &\\
  $ s = (e^{an})_{n \in \NN}, a \in \RR $ & $ \dfrac{z}{z - e^a} $ \\
        &\\
  $ s = (\sin(\omega n))_{n \in \NN}, \omega \in \RR $ & $ \dfrac{z \sin \omega}{z^2 - 2z \cos \omega + 1} $ \\
        &\\
  $ s = (\cos(\omega n))_{n \in \NN}, \omega \in \RR $ & $ \dfrac{z(z - \cos \omega)}{z^2 - 2z \cos \omega + 1} $
\end{tabular}

%%%%%%%%%%%%%%%%%%%%%%%%%%%%%%%%%%%%%%%%%%%%%%%%%%%%%%%%%%%%%%%%%%%%%% 

\subsection{Exerciții}

\indent\indent 1. Să se determine semnalul $x \in S_d^+$, a cărui transformată $Z$ este 
dată de:
\begin{enumerate}[(a)]
\item $\Lap_s (z) = \dfrac{z}{(z-3)^2}$ ;
\item $\Lap_s (z) = \dfrac{z}{(z-1)(z^2 + 1)}$ ;
\item $\Lap_s (z) = \dfrac{z}{(z-1)^2(z^2 + z - 6)}$ ;
\item $\Lap_s (z) = \dfrac{z}{z^2 + 2az + 2a^2}, a > 0$ parametru.
\end{enumerate}

\emph{Soluție:}

(a) Avem:
\begin{align*}
  x_n &= \dfrac{1}{2\pi i} \int_{|z| = \rho} z^{n-1} \Lap_s (z) dz \\
      &= \Rez(z^{n-1} \Lap_s(z), 3) \\
      &= \Rez \Big( \dfrac{z^n}{(z-3)^2}, 3 \Big) \\
      &= \lim_{z \to 3} \Big( (z-3)^2 \cdot \dfrac{z^n}{(z-3)^2} \Big)^{\prime} \\
      &= \lim_{z \to 3} nz^{n-1} \\
      &= n3^{n-1}.
\end{align*}

(b)
\begin{align*}
  x_n &= \dfrac{1}{2 \pi i} \int_{|z| = \rho} z^{n-1} \Lap_s(z) dz \\
      &= \Rez(z^{n-1} \Lap_s (z), 1) + \Rez(z^{n-1} \Lap_s(z), i) + 
        \Rez(z^{n-1} \Lap_s(z), -i) \\
  \Rez(z^{n-1}\Lap_s(z),1) &= \lim_{z \to 1} z^{n-1} \dfrac{z}{(z - 1)(z^2 + 1)} 
                             \cdot (z-1) = \dfrac{1}{2} \\
  \Rez(z^{n-1} \Lap_s(z), i) &= \dfrac{i^n}{2i \cdot (i-1)} \\
  \Rez(z^{n-1} \Lap_s(z), -i) &= \dfrac{(-1)^n i^n}{2i(i+1)}.
\end{align*}

Remarcăm că pentru $n = 4k$ și $n = 4k+1$, avem $x_n = 0$, iar în celelalte 
două cazuri, $x_n = 1$. \\

(c)
\begin{align*}
  \Rez(z^{n-1} \Lap_s(z), 1) &= \lim_{z\to 1} \Big[ z^{n-1} \cdot (z-1)^2 \cdot 
                               \dfrac{z^2}{(z-1)^2 \cdot (z^2 + z - 6)} \Big]^{\prime}\\
                             &= - \dfrac{4n + 3}{16}. \\
  \Rez(z^{n-1} \Lap_s(z), 2) &= \dfrac{2^n}{5} \\
  \Rez(z^{n-1} \Lap_s(z), -3) &= -\dfrac{(-3)^n}{80}.
\end{align*}

Obținem $x_n = - \dfrac{4n+3}{16} + \dfrac{2^n}{5} - \dfrac{(-3)^n}{80}$. \\

(d) $z_{1,2} = a (-1 \pm i)$ sînt poli simpli. Avem:
\begin{align*}
  x_n &= \Rez\Big( \dfrac{z^n}{(z^2 + 2a + 2a^2)}, z_1 \Big) + 
        \Rez\Big(  \dfrac{z^n}{(z^2 + 2a + 2a^2)}, z_2 \Big) \\
      &= \dfrac{a^n(-1 + i)^n}{2z_1 + 2a} + \dfrac{a^n(-1-i)^n}{2z_1 + 2a} \\
      &= -\dfrac{i}{2a}(z_1^n - z_2^n).
\end{align*}

Putem scrie trigonometric numerele $z_1$ și $z_2$:
\begin{align*}
  z_1 &= a(-1 + i) = a \sqrt{2} \Big( \cos \dfrac{3\pi}{4} + 
        i \sin \dfrac{3 \pi}{4} \Big) \\
  z_2 &= a(-1 - i) = a \sqrt{2} \Big( \cos \dfrac{3 \pi}{4} - 
        i \sin \dfrac{3 \pi}{4} \Big).
\end{align*}

Deci: $x_n = 2^{\frac{n}{2}} a^{n-1} \sin \dfrac{3 n \pi}{4}$.


%%%%%%%%%%%%%%%%%%%%%%%%%%%%%%%%%%%%%%%%%%%%%%%%%%%%%%%%%%%%%%%%%%%%%%%%%%%%%%%% 

\vspace{1cm}

2. Fie $x = (x_n) \in S_d^+$ și $y = (y_n)$, unde $y_n = x_0 + \dots + x_n$. 
Să se arate că $Y(z) = \dfrac{z}{z-1} X(z)$.

\emph{Soluție:}

Avem $Y(z) = \displaystyle\sum_{n=0}^\infty y_n z^{-n}.$ Dar: 
\[
  X(z) = \sum_{n=0}^\infty x_n z^{-n} \text{ și } 
  \sum_{n=0}^\infty x_{n-1} z^n = \dfrac{1}{z} X(z),
\]
deoarece $x_{-1} = 0$. Putem continua și obținem 
$\displaystyle\sum_{n=0}^\infty x_{n-k} z^{-n} = \dfrac{1}{z^k} X(z)$. Așadar: 
\[
  Y(z) = X(z) \cdot \Big( 1 + \dfrac{1}{z} + \dfrac{1}{z^2} + \dots \Big) 
  = X(z) \cdot \dfrac{z}{z-1}.
\]


\vspace{1cm}

3. Cu ajutorul transformării $Z$, să se determine șirurile $(x_n)$ definite 
prin următoarele relații:
\begin{enumerate}[(a)]
\item $x_0 = 0, x_1 = 1, x_{n+2} = x_{n+1} + x_n, n \in \mathbb{N}$ 
  (șirul lui Fibonacci);
\item $x_0 = 0, x_1 = 1, x_{n+2} = x_{n+1} - x_n, n \in \mathbb{N}$;
\item $x_0 = x_1 = 0, x_2 = -1, x_3 = 0, x_{n+4} + 
  2x_{n+3} + 3 x_{n+2} + 2 x_{n+1} + x_n = 0, n \in \mathbb{N}$;
\item $x_0 = 2, x_{n+1} + 3x_n = 1, n \in \mathbb{Z}$;
\item $x_0 = 0, x_1 = 1, x_{n+2} - 4 x_{n+1} + 3x_n 
  = (n+1)4^n, n \in \mathbb{N}$.
\end{enumerate}

\emph{Soluție:}

Abordarea generală este să considerăm șirul $(x_n)$ ca fiind restricția unui 
semnal $x \in S_d^+$ la $\mathbb{N}$ și rescriem relațiile de recurență sub 
forma unor ecuații de convoluție $a \ast x = y$, pe care le rezolvăm în $S_d^+$.\\

(a) Fie $x \in S_d^+$, astfel încît restricția lui la $\mathbb{N}$ să fie șirul 
căutat. Deoarece avem: 
\[
  x_{n+2} - x_{n+1} - x_n = y_n, n \in \mathbb{Z},
\]
cu $y_n = 0$ pentru $n \neq -1$ și $y_{-1} = 1$, avem ecuația de convoluție: 
\[
  a \ast x = y, \text{ unde } a 
  = \delta_{-2} + \delta_{-1} + \delta, y = \delta_{-1}.
\]

Aplicăm transformata $Z$ și rezultă: 
\[
  \Lap_s x(z) (z^2 - z - 1) = z \Longrightarrow x_n 
  = \dfrac{1}{\sqrt{5}} \Big[ \Big( \dfrac{1 + \sqrt{5}}{2} \Big)^n - 
  \Big( \dfrac{1 - \sqrt{5}}{2} \Big)^n \Big].
\]

(b) Ca în cazul anterior, avem $a \ast x = y$, cu 
$a = \delta_{-2} - \delta_{-1} + \delta$, unde $y = \delta_{-1}$. Aplicînd 
transformata $Z$, obținem: 
\[
  \Lap_s x(z) (z^2 - z + 1) = z \Longrightarrow \Lap_s x(z) 
  = \dfrac{z}{z^2 - z + 1}.
\]

Obținem:
\[
  x_n = \Rez \Big( z^{n-1} \dfrac{z}{z^2 - z + 1}, 
  \dfrac{1 + i \sqrt{3}}{2} \Big) + \Rez \Big( z^{n-1} \dfrac{z}{z^2 - z + 1}, 
  \dfrac{1 - i \sqrt{3}}{2} \Big).
\]

Calculăm reziduurile, cu notația $\varepsilon = \dfrac{1 + i \sqrt{3}}{2}$ și 
$\overline{\varepsilon} = \dfrac{1 - i \sqrt{3}}{2}$:
\begin{align*}
  \Rez \Big( \dfrac{z^n}{z^2 - z + 1}, \varepsilon \Big) 
  &= \lim_{z \to \varepsilon} \dfrac{z^n}{z^2 - z + 1} 
    \big( z - \varepsilon \big) \\
  &= \dfrac{\varepsilon^n}{i \sqrt{3}} \\
  &= \dfrac{\cos \frac{2 n \pi}{3} + i \sin \frac{2 n \pi}{3}}{i \sqrt{3}}.
\end{align*}

Similar: 
\[
  \Rez \Big( \dfrac{z^n}{z^2 - z + 1}, \overline{\varepsilon} \Big) 
  = \dfrac{\cos \frac{2n\pi}{3} - i \sin \frac{2 n \pi}{3}}{-i\sqrt{3}}.
\]

Rezultă: 
\[
  x_n = \dfrac{2}{\sqrt{3}} \sin \dfrac{2n\pi}{3}, n \in \mathbb{N}.
\]

(c) Ecuația $a \ast x = y$ este valabilă pentru: 
\[
  a = \delta_{-4} + 2 \delta_{-3} + 3 \delta_{-2} + 2 \delta_{-1} + \delta, 
  \quad y = -\delta_{-2} - 2 \delta_{-1}.
\]

Aplicăm transformata $Z$ și obținem: 
$\Lap_s x(z) = - \dfrac{z(z+2)}{(z^2 + z + 1)^2}$. Descompunem în fracții 
simple, calculăm reziduurile și ținem cont de faptul că rădăcinile numitorului, 
$\varepsilon_1, \varepsilon_2$ sînt poli de ordinul 2, obținem: 
\[
  x_n = \dfrac{(2n-4)(\varepsilon_1^n - \varepsilon_2^n) - 
    (n+1)(\varepsilon_1^{n-1} + \varepsilon_1^{n-2} - \varepsilon_2^{n-1} - 
    \varepsilon_2^{n-2})}{(\varepsilon_2 - \varepsilon_1)^3} 
  = \dfrac{2(n-1)}{\sqrt{3}} \sin \dfrac{2 n \pi}{3}, n \in \mathbb{N}.
\]

(d) Ecuația corespunzătoare este $a \ast x = y$, cu $a = \delta_{-1} + 3 \delta$ și 
$y_n = 1, \forall n \geq 1$, iar $y_{-1} = x_0 + 3 x_{-1} = 2$, cu 
$y_n = 0, \forall n \leq -2$, adică $y = 1 + 2 \delta_{-1}$.

Așadar: 
\[
  \delta_{-1} \ast x + 3 \delta \ast x = 1 + 2 \delta_{-1}.
\]

Aplicăm transformata $Z$ și obținem:
\begin{align*}
  z \Lap_s x(z) + 3 \Lap_s x(z) &= \dfrac{z}{z-1} + 2z \\
                                &= \dfrac{2z^2 + 3z}{z-1} \\
  \Longrightarrow \Lap_s x(z) &= \dfrac{2z^2 + 3z}{(z-1)(z+3)} \\
  \Longrightarrow x_n &= \Rez(z^{n-1} \cdot \dfrac{2z^2 + 3z}{(z-1)(z+3)}, 1) + 
                        \Rez (z^{n-1} \cdot \dfrac{2z^2 + 3z}{(z-1)(z+3)}, -3) \\
  \Rez(z^{n-1} \cdot \dfrac{2z^2 + 3z}{(z-1)(z+3)}, 1) &= \lim_{z \to 1} (z-1) 
                                                         \cdot z^{n-1} \cdot \dfrac{2z^2 + 3z}{(z-1)(z+3)} = \dfrac{5}{4} \\
  \Rez(z^{n-1} \cdot \dfrac{2z^2 + 3z}{(z-1)(z+3)}, -3) 
                                &= \lim_{z \to -3} (z+3) z^{n-1} \cdot \dfrac{2z^2 + 3z}{(z-1)(z+3)} \\
                                &= (-3)^{n-1} \cdot \dfrac{12}{-4} = -3 \cdot (-3)^{n-1}.	                              
\end{align*}

Rezultă: $x_n = \dfrac{5}{4} - 3 \cdot (-3)^{n-1}$. \\

(e) Avem ecuația: $a \ast x = y$, unde 
$a = \delta_{-2} - 4 \delta_{-1} + 3 \delta$, cu 
$y_n = 0, \forall n \leq -2, y_{-1} = 1$ și 
$y_n = (n+1)4^n, \forall n \in \mathbb{N}$.

Fie $s_1 = (n4^n)_n, s_2 = (4^n)_n$. Atunci:
\begin{align*}
  \Lap_s s_1(z) &= -z \Lap_s s_2^{\prime}(z) 
                  = - z \big( \dfrac{z}{z-4} \big)^{\prime} = \dfrac{4z}{(z-4)^2} \\
  \Longrightarrow \Lap_s x(z)(z^2 - 4z + 3) 
                &= \dfrac{4z}{(z-4)^2} + \dfrac{z}{z-4} + z \\
                &= \dfrac{z^2}{(z-4)^2} + z \\
  \Longrightarrow  \Lap_s x(z) &= \dfrac{z(z^2 - 7z + 16)}{(z-4)^2 (z-1) (z-3)}.
\end{align*}

Descompunem în fracții simple și obținem, în fine: 
\[
  x_n = \dfrac{1}{9} \big[ 18 \cdot 3^n + (3n - 13) 4^n - 5 \big], 
  n \in \mathbb{N}.
\]

\textbf{OBSERVAȚIE:} Toate exercițiile cu recurențe se mai pot rezolva în alte
două moduri:

(1) Se poate aplica teorema de convoluție relației de recurență. De exemplu,
din recurența:
\[
  x_{n + 2} - 2x_{n + 1} + x_n = 2
\]
putem obține:
\[
  Z(x_{n+2}) - 2Z(x_{n + 1}) + Z(x_n) = 2Z(1),
\]
iar $ Z(x_{n+2}) = Z(x_n \ast \delta_{-2}) = Z(x_n) \cdot Z(\delta_{-2}) $ etc.

(2) Se poate aplica teorema de deplasare. În aceeași recurență de mai sus,
de exemplu, avem:
\[
  Z(x_{n+2}) = z^n \Big( Z(x_n) - x_0 - x_1 z^{-1} \Big)
\]
și la fel pentru celelalte.


%%% Local Variables:
%%% mode: latex
%%% TeX-master: "../m3-19-20"
%%% End:
